
\documentclass[../../../physics-standard_questions_all.tex]{subfiles}


\begin{document}

図のように，真空中において，同形の3枚の金属極板X，Y，Zを
平行に狭い間隔で向き合わせてあり，電池B$_1$（起電力$2V_0$），B$_2$（起電力$V_0$），
スイッチS$_1$，S$_2$およびアースと接続されて回路を形成している．
XとYの間隔は$d$，YとZの間隔は$2d$である．
極板の面積を$S$，真空の誘電率を$\varepsilon_0$とし，端の効果は無視する．
はじめ，S$_1$，S$_2$は共に開いており，いずれの極板も帯電していない．

\begin{enumerate}

    \item S$_1$を閉じて安定するまで待った．
    このとき，X下面の電荷$x_1$，Y下面の電荷$y_1$を求めよ．

    \item S$_1$を開き，S$_2$を閉じて安定するまで待った．
    このとき，X下面の電荷$x_2$，Y下面の電荷$y_2$を求めよ．
   
    \item S$_2$を開き，S$_1$を閉じて安定するまで待った．
    このとき，X下面の電荷$x_3$，Y下面の電荷$y_3$を求めよ．
    
    \item スイッチはそのままで，XとYの間を誘電体（比誘電率$\varepsilon_r$）で満たし，
    安定するまで待った．このとき，X下面の電荷$x_4$，Y下面の電荷$y_4$を求めよ．

\end{enumerate}

\vspace{0.2\baselineskip}
{\centering
    \includegraphics{\subfix{fig/fig_em_cb_11/fig_em_cb_11_00.pdf}}
\par}
\vspace{0.2\baselineskip}

\end{document}



